We denote by $\HH_g$ the {\it Siegel upper half-space}~-- the
space of complex symmetric $g\times g$ matrices with positive definite
imaginary part. An element $\tau\in\HH_g$ is called a {\it period
matrix} and defines the complex abelian variety $X_\tau:=\CC^g/\ZZ^g+\tau
\ZZ^g$. It is a well-known fact that the moduli space of principally polarized abelian varieties (ppavs for short) can be identified with the quotient of $\HH_g$ by an action of the symplectic group ${\rm Sp}(2g, \ZZ)$ which is a generalization of the standard action of ${\rm SL}(2,\ZZ)$ on the complex upper-half plane, namely,
\[
\gamma \cdot \tau := (a \tau + b) \cdot (c \tau + d)^{-1},\qquad \forall\, \gamma = \begin{pmatrix} a & b \\ c & d \end{pmatrix} \in {\rm Sp}(2g,\ZZ),
\]
where $a$, $b$, $c$, $d$ are $g \times g$ blocks. A ppav is called {\it irreducible} if it is not isomorphic to a product of two lower-dimensional ppavs.

For $\ep,\de\in \ZZ_2^g$ and $z\in \CC^g$ we define the {\it first order theta function with characteristic $[\ep,\de]$} to be
\[
\tt\ep\de(\tau,z):=\sum\limits_{m\in\ZZ^g} \exp \pi {\rm i} \left[\left(
m+\frac{\ep}{2}\right)^t\tau \left(m+\frac{\ep}{2}\right)+2\left(m+\frac{\ep}{2}\right)^t\left(z+
\frac{\de}{2}\right)\right].
\]

Characteristics can be defined for $\ep$, $\de$ in $\ZZ^g$, but the {\it reduction formula}
\begin{gather*}%\label{eq:Ad}
\tt{\ep+2\ep'}{\de+2\de'} (\tau , z) = (-1)^{\langle \ep ,\de'\rangle} \tt\ep\de(\tau,z) \end{gather*}
(where the symbol $\langle \cdot , \cdot \rangle$ stands for the standard inner product) shows that these functions are uniquely determined up to a sign by considering~$\ep$ and~$\de$ as vectors of zeros and ones. Henceforward we shall work with reduced characteristics with the agreement that a theta function associated with a sum of characteristics is meant to be non-reduced.

Characteristics are defined {\it even} or {\it odd} depending on whether respectively $\langle \ep , \de\rangle =0,\, =1$ ${\rm mod}\, 2$. A straightforward computation shows there are $ 2^{g-1}\big( 2^{g}+1\big)$ even characteristics and $2^{g-1}\big( 2^{g}-1\big)$ odd ones. It is easily seen that theta functions associated with even characteristics are even functions in the variable~$z$, whereas those associated with odd characteristics are odd functions.
Triplets of characteristics $[\ep,\de]$, $[\ep',\de']$, $[\ep'',\de'']$ are also called {\it azygetic} if $\langle \ep , \de\rangle+\langle\ep' , \de'\rangle+\langle\ep'' , \de''\rangle+\langle\ep + \ep' + \ep'' , \de + \de' + \de''\rangle=1$. More generally, a $k$-tuple of characteristics is called azygetic if its subtriplets are azygetic.

Theta functions satisfy an addition formula (cf.~\cite{Ig72} for a general formulation and~\cite{Gr} or~\cite{GSM} for the specific version we use here)
\[
\theta \begin{bmatrix} \ep \\ \de \end{bmatrix} (2\tau, 2z) \theta \begin{bmatrix} \ep + \ep' \\ \de \end{bmatrix} (2\tau, 2w)= \frac{1}{2^g} \sum_{\s \in \ZZ_2^{g}}(-1)^{\langle \ep , \s\rangle} \theta \begin{bmatrix} \ep '\\ \de +\s \end{bmatrix} (\tau, z+w) \theta \begin{bmatrix} \ep'\\ \s \end{bmatrix} (\tau,z-w).
\]

For a given $\tau \in H_g$ we denote by $X_{\tau}[2]$ the group of points of order two on $X_{\tau}$ and by $L_{\tau}$ a~symmetric line bundle on $X_\tau$ defining the principal polarization. Note that
\[
\tt\ep\de(\tau,z)=\exp \pi {\rm i} \left[\frac{\ep^t}{2}\tau\frac{\ep}{2} +\frac{\ep}{2}\left(z+\frac{\de}{2}\right) \right]\tt{0}{0}\left( \tau,z+ \tau\frac{\ep}{2} +\frac{\de}{2} \right)
\]
and the function $z \to
\theta\left[\begin{smallmatrix}\ep \\ \de \end{smallmatrix}\right] (\tau,z)$ defines the unique (up to scalar multiplication)
section of $t_x^*L_{\tau}$, i.e., the translate of the line bundle $L_{\tau}$ by the point of order two $x=(\tau \ep +\de)/2$. Because of the duplication map, the functions $z \rightarrow
\theta\left[\begin{smallmatrix}\ep \\ \de \end{smallmatrix}\right] (\tau,2z)$ are a basis of $H^0\big(X_\tau, L_{\tau}^{\otimes 4}\big)$.

For $\ep\in\ZZ_2^g$ we also
define the {\it second order theta function with characteristic $\ep$} to
be
\[
\T[\ep](\tau,z):=\tt{\ep}{0}(2\tau,2z).
\]
The functions $z \rightarrow \T[\ep](\tau,z)$ are a basis of $H^0\big(X_\tau, L_{\tau}^{\otimes 2}\big)$. Since these functions are even, i.e., $
\T[\ep](\tau,-z)=
\T[\ep](\tau,z)$, they induce a map
\[
{\rm Th}_2\colon \ X_\tau\to \PP^{2^g-1}
\]
 that factorizes along the Kummer variety $ K_{\tau}:=X_\tau/\pm 1$.

By evaluating first order theta functions at $z=0$ we get the so-called {\it theta constants} $\theta\left[\begin{smallmatrix}\ep \\ \de \end{smallmatrix}\right] (\tau,0)$; because of the parity of theta functions, theta constants associated with odd characteristics are clearly trivial. Analogously, second order theta functions evaluated at $z=0$ yield {\it second order theta constants} $\T[\ep](\tau,0)$. As functions of $\tau$ these are not well defined on the moduli space of ppavs but, as a consequence of how they transform under the action of ${\rm Sp}(2g, \ZZ)$, they induce maps on finite covers of the moduli space. More precisely, if we introduce the following family of subgroups of finite index in ${\rm Sp}(2g, \ZZ)$
\begin{gather*}
\Gamma_g[n]:= \operatorname{Ker}({\rm Sp}(2g,\ZZ) \to {\rm Sp}(2g,\ZZ_n)), \qquad n \in \NN,\\
\Gamma_g[n,2n]:=\left\lbrace\gamma = \begin{pmatrix} a & b \\ c & d \end{pmatrix} \in \Gamma_g[n] \,|\, {\rm diag}\big(a^tb\big)\equiv{\rm diag}
\big(c^td\big)\equiv0\ {\rm mod}\ 2n \right\rbrace,
\end{gather*}
first order theta constants are seen to transform as follows (cf.~\cite{Ig72} and~\cite{RF} for the general formula)
\[
\tt\ep\de(\gamma \cdot \tau,0) = \kappa(\gamma) \chi_{\ep,\de}(\gamma) \det{(c \tau + d)}^{\frac{1}{2}} \tt\ep\de(\tau,0) \qquad \forall\, \gamma \in \Gamma_g[2],
\]
where $\kappa(\gamma)$ is an eighth root of the unity for any $\gamma$ and $ \chi_{\ep,\de}$ are characters of the group $\Gamma_g[2,4]/\Gamma_g[4,8]$; in particular, an action of $\Gamma_g[2,4]/\Gamma_g[4,8]$ is naturally induced on theta constants and can be described in terms of these characters $ \chi_{\ep,\de}$, which span the character group of $\Gamma_g[2,4]/\pm\Gamma_g[4,8]$ (cf.~\cite{SM}). Thanks to the transformation formula, first order theta constants induce a map of $\HH_g/\pm\Gamma_g[4,8]$ into the projective space; this map is known to be an immersion, cf.~\cite{Ig72}. Second order theta constants also induce a map on a finite cover of the moduli space of ppavs
\[
\mathcal{T}\colon \ \HH_g /\Gamma_g[2,4] \to \PP^{2^g-1},
\]
which is generically injective for any $g$, cf.~\cite{SM}. Once a basis $\omega_1, \dots , \omega_g$ for the cohomology of an algebraic curve of genus $g$ is chosen together with a symplectic basis of cycles $\de_1, \dots , \de_g, \de'_1, \dots, \de'_g$, the $g \times 2g$ matrix $\big(\int_{\de_j}\omega_i, \int_{\de'_j}\omega_i\big)$ defines a complex torus which is isomorphic to a ppav $X_{\tau}$; this complex torus is known as the {\it Jacobian variety} of the curve and the locus in the moduli space of ppavs defined by those ppavs that are isomorphic to Jacobians of curves is known as the {\it Jacobian locus}. Because of the existence of the map $\mathcal{T}$, one knows that the Jacobian locus and the {\it hyperelliptic locus} (i.e., the locus defined by those ppavs that are isomorphic to Jacobians of hyperelliptic curves) can be described in terms of equations involving theta constants by taking the preimages of these loci under the covering map $\HH_g /\Gamma_g[2,4] \to \HH_g /{\rm Sp}(2g, \ZZ)$; the preimage of the hyperelliptic locus has several irreducible components in $\HH_g /\Gamma_g[2,4]$.

As for the hyperelliptic locus, we know from the result in~\cite{mu} that the irreducible components are defined in terms of vanishing and non-vanishing conditions for certain theta constants $\theta\left[\begin{smallmatrix}\ep \\ \de \end{smallmatrix}\right] (\tau,0)$. Mumford's result is based on Neumann's dynamical system. In~\cite{po} Poor used the cross ratio to prove that the vanishing conditions alone are sufficient to characterize the hyperellipitic locus, once the irreducibility of the principally polarized abelian variety is assumed. In both cases a fundamental tool is the so-called Frobenius theta formula, which is a consequence of Riemann's theta formula (cf.~\cite{FR} and~\cite{mu}). A~few years ago in \cite{Gr}, a characterization of the hyperelliptic abelian varieties was given in terms of cubic equations in the second order theta functions. These cubics were obtained by using the explicit coefficients for the addition formula from~\cite{BK2}.

The aim of this note is twofold: we first relate Grushevsky's approach to the others' by proving that his cubic equations are related to Frobenius' theta formula; we then give a new proof of Mumford's characterization by applying Gunning's multisecant formula (cf.~\cite{Gu}), which is a remarkable geometrical characterization of the locus of Jacobians in terms of intersections of g-dimensional linear subspaces with the Kummer variety of an abelian variety.

In this way we have an explicit relation between various approaches; in particular, we get an immediate link between the geometrical characterization and the explicit equations for the hyperelliptic locus.

Incidentally, we also prove that Grushevsky's cubics are induced by the quadratic relations related to the theta vanishing.

\section{Theta functions}

Throughout the rest of the paper we shall omit the subscript $\tau$ whenever it is clear that $\tau$ is fixed.
For the topics introduced here we shall follow Section 1 in \cite{vG} closely. Once we set $ M:= L^{\otimes 2}$, as $t_{x}^*M\cong M$ if and only if $x \in X[2]$, the following group is well defined
\[
G(M):=\big\{(x, \phi) \,|\, x\in X[2], \, \phi\colon t_{x}^*M \xrightarrow{\cong} M \big\}
\]
and an irreducible action of $G(M)$ on the space of global sections of $M$ is also defined by setting $(x, \phi) s := \phi (t_{x}^*s)$.

A theta structure is an isomorphism between $G(M)$
and the Heisenberg group, namely the set
\[
H:=\CC^{*}\times\ZZ_{2}^g\times {\rm Hom}(\ZZ_{2}, \CC^{*})^g,
\]
provided with the group law
\[
(t,x,x^*) \cdot (s,y,y^*) :=(tsy^*(x), x+y, x^*+y^*).
\]

Now, let $B_n:=\Gamma(X, M^n)$ for any $n \geq 1$.
The action of $\{ \pm 1 \}\subset {\rm End}( X)$ decomposes each of these spaces into the direct sum of two factors $B_n ^{+}$ and $ B_n ^-$. One has $B_1= B_1^+$ and $B_1$ is an irreducible representation of $H$.
Moreover, $B_1$ admits a basis $\{X_{\s}\}$ with $\s\in\ZZ_2^g$ such that
$(t, x, x^*)X_{\s}=tx^*(x+\s)X_{x+\s}$. The Heisenberg group naturally acts on the $n$-fold symmetric products $S^nB_1$ as well. To decompose~$S^nB_1$ we introduce the two following subgroups of $H$:
\[
K=\{(t, x, x^*)\colon t=1,\, x=0\} , \qquad K^*=\{(t, x, x^*)\colon t=1,\, x^*=0\},
\]
and we recall from \cite{vG} the following:

\begin{Proposition}\label{pr.vG1}\quad
\begin{enumerate}\itemsep=0pt
\item[$(i)$] A basis of eigenvectors for the action of H on $S^2B_1$ is given by the $2^{g-1}\big( 2^{g}+1\big)$ elements $Q[\ep, \ep']=\sum(-1)^{\langle \s,\ep'\rangle}X_{\s}X_{\s+\ep}$, where $\s, \ep, \ep'\in\ZZ_2^g$, the symbol $\langle \cdot,\cdot\rangle$ stands for the standard inner product and $\langle \ep,\ep'\rangle=0$. The action of $H$ on the elements of the basis is given by
\[
(t, x, x^*)Q[\ep, \ep']=t^2(-1)^{\langle x,\ep'\rangle}x^*(\ep)Q[\ep, \ep'].
\]
\item[$(ii)$] The space $S^3B_1$ is a direct sum of $\big(2^g+1\big)\big(2^{g-1}+1\big)/3$ copies of a~$2^g$-dimensional irreducible representation of~$H$.
\item[$(iii)$]Let $X(H)$ be the group of characters of $\ZZ_{2}^g\times {\rm Hom}(\ZZ_{2}, \CC^{*})^g$. Then
\[
S^4B_1= \bigoplus_{\chi\in X(H)} S^4_{\chi}(B_1).
\]
The dimension of the eigenspace associated with the trivial character is equal to $\big(2^g+1\big)\big(2^{g-1}+1\big)/3$, whereas the dimension of the other eigenspaces is $\big(2^{g-1}+1\big)\big(2^{g-2}+1\big)/3$.
\end{enumerate}
\end{Proposition}

Moreover, a basis for $S^3B_1$ as an $H$-module is given by a basis of the $K$-invariant elements and a basis for $S^4_{0}(B_1)$ is also described in \cite{vG}.

Maps between these spaces can be defined. Let $S^3_{0}(B_1)$ be the subspace of $K$-invariant elements of $S^3(B_1)$.
 For $F\in S^3_{0}(B_1)$ and $ \s\in \ZZ_{2}^g$ let
\[
F_{\s}:=( 1, \s, 0)F.
\]
For $\chi=(y^*, y)\in X(H)$ (so that $\chi(x, x^*)=y^*(x)x^*(y)$) we define
\begin{gather*}
M(\chi)\colon \ S^3_{0}(B_1)\to S^4_{\chi}(B_1),\\
M(\chi)(F)= \sum_{\s\in \ZZ_{2}^g} y^*(\s)X_{\s+y}F_{\s}.
\end{gather*}

 \begin{Proposition} \label{pr.vG2}The maps $M(\chi)$ are surjective. In the case $\chi=0$ the map $M(0)$ is an isomorphism and the inverse is given by $ 4^{-1}\frac{\partial}{\partial X_0}$. Moreover
\[
\left(\frac{\partial}{\partial X_{\s}}\right)(M(0)F)=4F_{\s},\qquad \forall\, \s\in \ZZ_{2}^g.
\]
\end{Proposition}

We will briefly discuss the case of the quartics; more details on the subject can be found in~\cite{vG}.

The decomposition of $S^4B_1$ is given in Proposition~\ref{pr.vG1}(iii); a similar decomposition holds for the space
\[
S^2\big(S^2B_1\big)=\bigoplus_{\chi} \big(S^2\big(S^2B_1\big)\big)_{\chi}.
\]

In this case the dimensions of the eigenspaces are respectively equal to $2^{g-1} \big(2^g+1\big)$ and $2^{g-2}\big(2^{g-1}+1\big)$. Obviously, we have a surjective map
\[
 \Phi_4\colon \ S^2\big(S^2B_1\big)\to S^4B_1.
\]
 According to \cite{fay} (cf.\ also~\cite{APS} or~\cite{FGS}), the relations are generated by a particular case of the so-called Riemann relations, namely those of the form
\begin{gather}\label{eq:RR}
\sum_{\langle\ep,\de\rangle =0} (-1)^{\langle \s, \de\rangle} v_{\ep,\de}Q[\ep, \de]Q[\ep+\s, \de+\r] .
\end{gather}

Here the vector $v=( v_{\ep,\de})$ varies in a suitable space of dimension $\big(2^{2g}-1\big)/3$. We say that these are biquadratic relations. By evaluating at the theta functions
$X_{\s}=\T[\s](\tau,z)$ and applying the addition formula, we get
\begin{gather}\label{eq:vT}
Q[\ep, \ep'](\dots,\T[\s](\tau,z),\dots):=Q[\ep, \ep'](\tau, z) =\tt \ep {\ep'} (\tau, 0)\tt \ep {\ep'} (\tau, 2z).
 \end{gather}
Hence, a basis for the quadrics containing the Kummer variety is given by those $Q[\ep, \ep']$ such that
$\theta\left[\begin{smallmatrix}\ep \\ \ep' \end{smallmatrix}\right] (\tau, 0)=0$.

To sum up, we can also state the following:

\begin{Proposition} \label{trivial relations}
All biquadratic Riemann relations induce trivial relations between second order theta functions.
\end{Proposition}

 We know there also exist highly non-trivial relations like the equation defining the Kummer surface in genus 2 or the Coble quartic in genus~3. These are in the kernel of the map
\[
\Psi_4\colon \ S^4B_1\to B_4^+.
\]
\begin{Remark} More general Riemann relations are obtained by evaluating (\ref{eq:RR}) at the theta functions and then at the points $z+(\tau x+y)/4$, $x,y \in\ZZ^g$, which yields
\begin{gather}\label{eq:RR2}
\sum_{\langle\ep,\de\rangle =0}\! (-1)^{\langle\s+x, \de\rangle } v_{\ep,\de} \tt\ep\de(\tau,0)\tt{\ep+\s }{\de+\r}(\tau,0)\tt{\ep+x}{\de+y}(\tau,2z)\tt{\ep+\s+x}{\de+\r+y}(\tau,2z)=0.\!\!\!
\end{gather}
\end{Remark}
